\subsection{Global section coordinates by actual norm systems}

\begin{statement}{E.46}{Proved}{Every global fiber for the inherited module quotient, and explicit ramified families}
Use the actual arithmetic quotient of E.38,
\[
 \Lambda=\mathbb Z_3[[U]],\quad T=1+U,\quad
 J=(3,U)\Lambda,\quad H=3\mathbb Z_3\ltimes J,
 \qquad
 g(a,v)=(\chi_g a,\chi_gv(T^{\chi_g})).
 \tag{E46.1}
\]
Fix a continuous first-coordinate cocycle
$A\in Z^1(G_{\mathbb Q},3\mathbb Z_3(1))$.
For $n\geq0$ let $X_{A,n}=\mathbb Z/3^n$ with the
actual action
\[
             g i=A_g+\chi_g i\pmod{3^n},
 \qquad
 E_{A,n}=\operatorname{Map}(X_{A,n},\overline{\mathbb Q})^{G_{\mathbb Q}}.
 \tag{E46.2}
\]
The action on functions is $(g f)(i)=g(f(g^{-1}i))$.
The reduction $X_{A,n+1}\to X_{A,n}$ induces an
inclusion of finite etale algebras $E_{A,n}\to E_{A,n+1}$
of rank three, with its usual multiplicative norm.
In particular $E_{A,0}=\mathbb Q$.
For a finite product of number fields $E$ write
\[
 \widehat{E^\times}_3
    =\varprojlim_m E^\times/(E^\times)^{3^m},
 \qquad
 \mathcal N_A=
  \left\{(u_n)_n:
    u_n\in\widehat{E_{A,n}^\times}_3,
    \ N_{E_{A,n+1}/E_{A,n}}u_{n+1}=u_n\right\}.
 \tag{E46.3}
\]
These completions carry the inverse-limit topology of
the discrete power-class groups. They are not assumed
compact. All the transition maps in (E46.3) are the
indicated norms, including when either algebra is a
product.

The global sections of $H\rtimes G_{\mathbb Q}$ whose
first coordinate has class $[A]$, modulo conjugation
by $H$, are exactly
\[
 \boxed{\quad
  \mathcal N_A^{(3)}=
  \{(u_n)_n\in\mathcal N_A:
             u_0\in(\widehat{\mathbb Q^\times}_3)^3\}.
 \quad}                                                    \tag{E46.4}
\]
This is an isomorphism of topological abelian groups
after taking the section $g\mapsto(A_g,0)$ as origin.
The proof gives the map and its inverse, the exact
change-of-origin maps, and the continuous-cochain
topology. No surjectivity of the norm maps is needed.
Every first-coordinate class occurs.

For a rational $a\in\mathbb Q^\times$, take
$A=3\kappa_3(a)$ with compatible Kummer roots. Then
\[
 E_{A,n}=\mathbb Q[X_n]/(X_n^{3^n}-a^3),
 \qquad X_n\longmapsto X_{n+1}^3.                    \tag{E46.5}
\]
For $a^3\ne1$ and every $r,z\in\mathbb Z_3$, the explicit
tuple
\[
                 u_{n,r,z}=X_n^r(1-X_n)^{3z}        \tag{E46.6}
\]
defines a section class in (E46.4). Its image under the
actual homomorphism
\[
 H\longrightarrow\mathbb Z_3(1)^2,
 \qquad(a,v)\longmapsto
       \bigl(a/3,\operatorname{aug}(v)/3\bigr)
 \tag{E46.7}
\]
is
\[
             \bigl(\kappa_3(a),
                       r\kappa_3(a)+z\kappa_3(1-a^3)\bigr).
 \tag{E46.8}
\]
The parameter $z$ distinguishes the subfamily $r=0$.
For $a=2$, both parameters are distinguished by the
valuations at two and seven, giving an embedded
$\mathbb Z_3^2$ family in this fiber. In particular
$a=2,r=0,z=1$ gives the rational toric Kummer pair
$(2,-7)$ in this quotient.

For this same first coordinate $a=2$, the entire
norm system in (E46.4), including its field factors,
has the following exact description. For $n\geq1$
choose
\[
 \alpha_n=2^{1/3^{n-1}},\qquad
 \beta_n=\alpha_n\zeta_{3^n},\qquad
 F_n=\mathbb Q(\alpha_n),\qquad L_n=\mathbb Q(\beta_n).
 \tag{E46.9}
\]
Then
\[
 \begin{gathered}
 E_{A,n}=F_n\times L_n,
 \quad [F_n:\mathbb Q]=3^{n-1},
 \quad [L_n:\mathbb Q]=2\cdot3^{n-1},\quad
 L_1=\mathbb Q(\zeta_3),\\
 \alpha_{n+1}^3=\alpha_n,\qquad
 \beta_{n+1}^3=\beta_n .
 \end{gathered}                                          \tag{E46.10}
\]
Thus (E46.4) consists precisely of two sequences
$v_n\in\widehat{F_n^\times}_3$ and
$w_n\in\widehat{L_n^\times}_3$, each compatible under
the displayed degree-three field norms, with
\[
               v_1 N_{L_1/\mathbb Q}(w_1)
                         \in(\widehat{\mathbb Q^\times}_3)^3.
 \tag{E46.11}
\]
Formula (E46.6) is the pair
$(\alpha_n^r(1-\alpha_n)^{3z},
\beta_n^r(1-\beta_n)^{3z})$ in these coordinates.

\emph{Manuscript: \S51.16, folios733v--736v,
scans172--178, Part E(c),(d). The arithmetic group here
is the quotient $H$ in E.38. The actual maps from the
intrinsic metabelian and marked groups have exactly
the kernels (E38.9); (E46.4) classifies their target
section set. It does not supply a lift through either
of those kernels.}
\end{statement}
\begin{proof}
\emph{The twisted module and geometric conjugation.}
Write a section cocycle as $(A_g,b_g)$. The multiplication
in (E46.1) gives
\[
 b_{gh}=b_g+\chi_gT^{A_g}b_h(T^{\chi_g}).
 \tag{E46.12}
\]
Define $\Lambda_A$ to be the compact additive group
$\Lambda$ with action
\[
                g\cdot_A f=\chi_gT^{A_g}f(T^{\chi_g}),
 \tag{E46.13}
\]
and let $J_A$ be its stable submodule $J$.
The cocycle identity for $A$ verifies the action law
in (E46.13). Thus sections with their first cocycle
equal to $A$ have precisely the second cocycles
$Z^1(G_{\mathbb Q},J_A)$.

For completeness, conjugating by an actual element
$(s,u)\in H$ gives
\[
 \begin{split}
 A'_g&=A_g+(1-\chi_g)s,\\
 b'_g&=T^s b_g+u-\chi_gT^{A'_g}u(T^{\chi_g}).
 \end{split}                                             \tag{E46.14}
\]
The first coordinate is unchanged exactly when
$s=0$, since complex conjugation has $\chi=-1$
and $3\mathbb Z_3$ is torsion-free. The remaining
identifications in (E46.12) are exactly the
$J_A$ coboundaries. Every cocycle whose first class
is $[A]$ can be brought to first cocycle $A$ by the
unique required $s$. This change is continuous:
evaluate the first-coordinate difference at complex
conjugation and divide by two. Consequently the
fiber in the statement is precisely
$H^1_{\rm cts}(G_{\mathbb Q},J_A)$, with its quotient
topology and its origin $b=0$.

Equation (E46.14) also supplies the exact comparison
between representatives of a first-coordinate class.
If $A'_g=A_g+(1-\chi_g)s$, the map
$i\mapsto i+s$ is an equivariant bijection
$X_{A,n}\to X_{A',n}$. Multiplication by $T^s$
intertwines $\Lambda_A$ and $\Lambda_{A'}$ and
induces these bijections at every positional level.
The corresponding algebra isomorphisms commute with
the norms. Thus the description of the fiber is
transported by explicit maps, rather than by an
unspecified normalization.

\emph{Finite positional modules and Shapiro's map.}
Set
\[
 M_{A,n}=\mathbb Z_3(1)\otimes_{\mathbb Z_3}
                   \mathbb Z_3[X_{A,n}],
 \qquad
       g(\delta_i)=\chi_g\delta_{A_g+\chi_gi}.
 \tag{E46.15}
\]
The coefficient transition sends $\delta_i$ to
$\delta_{i\bmod3^n}$. The actual group-ring
identification is
\[
 \Lambda_A=\varprojlim_n M_{A,n},
 \qquad
 T^i\longmapsto\delta_{i\bmod3^n}.                    \tag{E46.16}
\]
Its topology is the inverse limit over both the
positional level and the coefficient modulus $3^m$,
equivalently the $(3,U)$-adic topology. The coefficient
projections are surjective and split continuously
as maps of $\mathbb Z_3$-modules; equivariance of a
chosen splitting is not required.

Decompose $X_{A,n}$ into orbits and choose a point
$i_0$ in each. Its stabilizer is $G_{F}$, where $F$
is the corresponding field factor of $E_{A,n}$.
For a cocycle $c:G_{\mathbb Q}\to M_{A,n}$, Shapiro's
map is explicitly
\[
                     h\longmapsto c(h)_{i_0}
                       \quad(h\in G_F).              \tag{E46.17}
\]
Here the subscript denotes the coefficient of
$\delta_{i_0}$. To verify the map and its inverse
directly, choose $g_i$ with $g_i i_0=i$ and
$g_{i_0}=1$. A continuous cocycle
$z:G_F\to\mathbb Z_3(1)$ induces
\[
 c_z(g)_i=
  \chi(g_i)\,
  z\bigl(g_i^{-1}g g_{g^{-1}i}\bigr).                 \tag{E46.18}
\]
The argument of $z$ fixes $i_0$. Its product identity
proves the cocycle equation, and (E46.17) recovers
$z$. Conversely, applying (E46.18) to (E46.17)
gives $c_z=c+\delta m$, where
$m_i=c(g_i)_i$ and $\delta m(g)=g m-m$.
These formulas are continuous
and show the precise degree-one Shapiro isomorphism.
They agree with \cite[II, Proposition1.11 and
Example1.27(a)]{E_MilneCFT}.

The Kummer sequence and Hilbert's theorem 90
\cite[II, Proposition1.22 and its proof]{E_MilneCFT} give
\[
 H^1(G_F,\mu_{3^m})=F^\times/(F^\times)^{3^m}.
\]
Passing to $\mathbb Z_3(1)$ gives
$H^1_{\rm cts}(G_F,\mathbb Z_3(1))=
\widehat{F^\times}_3$. Indeed the finite groups
$\mu_{3^m}(F)$ satisfy the Mittag--Leffler condition,
so their first inverse-limit derivative is zero.
The continuous-cochain comparison is
\cite[Proposition1.6 and Theorem2.2]{E_JannsenContinuous}.
The same conclusion follows by aligning the finite
Kummer cocycles and taking their continuous limit.
Thus (E46.17) supplies actual isomorphisms
\[
       H^1_{\rm cts}(G_{\mathbb Q},M_{A,n})
                       \simeq\widehat{E_{A,n}^\times}_3.
 \tag{E46.19}
\]

Under these maps the positional pushforward is
corestriction, and therefore the multiplicative
norm. Here is a direct check of this identification.
For a field inclusion $F\subset L$, choose a
$3^m$-th root of $u\in L^\times$ and multiply its
conjugates over the embeddings of $L$ into
$\overline{\mathbb Q}$ over $F$. Their product is
a $3^m$-th root of $N_{L/F}u$. The Galois ratios of
this product are exactly the sum of the transferred
Kummer cocycles, with the coefficient twists in
(E46.18). A different choice of coset representatives
changes the root by the corresponding coboundary.
For products of fields, multiply the norms of the
factors. This proves the assertion at every finite
coefficient level and hence after completion.
It is the explicit transfer construction of
\cite[II, Example1.29]{E_MilneCFT}.

The direction of the residual deck action can
also be read from these maps. When $A=0$, write
$D_\gamma(\delta_i)=\delta_{\gamma i}$ on a unit
orbit, and choose $g$ with $\chi_g=\gamma$.
For $h$ fixing the chosen unit root, the cocycle
identity gives
\[
 \chi_g c(g^{-1}hg)_1
       =c(h)_\gamma+(\chi_h-1)c(g)_\gamma.
\]
The left side is the usual Galois action on its
Kummer cohomology class. The right side is
Shapiro evaluation of $D_{\gamma^{-1}}c$, up
to the displayed coboundary. Consequently the
usual field action corresponds to inverse deck
translation. This explicitly fixes the convention
for the cyclotomic norm system in E.47.

\emph{The positional inverse limit without a hidden term.}
For every number field $F$,
$H^0(G_F,\mathbb Z_3(1))=0$: its group of roots of
unity is finite, so it contains no compatible
nonzero Tate-module element. By the degree-zero
version of (E46.17),
\[
                    H^0(G_{\mathbb Q},M_{A,n})=0
                         \quad\hbox{for every }n.     \tag{E46.20}
\]
We now prove the needed inverse-limit statement
on actual cocycles. Given compatible classes
$[c_n]\in H^1(G_{\mathbb Q},M_{A,n})$, choose a
representative $c_0$. Inductively, the projection
of a representative at level $n+1$ differs from
$c_n$ by a coboundary $\delta b_n$. Lift $b_n$
through the surjective coefficient projection and
subtract its coboundary. This makes the cocycles
compatible. Their limit is continuous by
(E46.16), and is a cocycle in $\Lambda_A$.

If a compatible cocycle is a coboundary at every
level, write $c_n=\delta b_n$. The differences
between the projection of $b_{n+1}$ and $b_n$
are invariant and hence zero by (E46.20).
Thus the $b_n$ themselves have a compatible limit,
which bounds the original cocycle. This proves
\[
 H^1_{\rm cts}(G_{\mathbb Q},\Lambda_A)
    \xrightarrow{\sim}
       \varprojlim_n H^1_{\rm cts}(G_{\mathbb Q},M_{A,n})
    \xrightarrow{\sim}\mathcal N_A.                  \tag{E46.21}
\]
In particular no first inverse-limit derivative of
degree-zero cohomology has been discarded, and the
proof makes no assertion that a norm map on
multiplicative classes is surjective.

These maps also prove the claimed topology. Equip
continuous cocycles with uniform convergence in
finite coefficient-and-positional quotients, then
take the quotient by continuous coboundaries.
If a cocycle becomes a coboundary in one such
finite quotient, its bounding element lifts to
$\Lambda_A$. Subtracting that lifted coboundary
gives a representative zero in that quotient.
Therefore the kernels of the finite cohomology
maps give exactly the quotient neighborhoods.
Under (E46.19) these are the power-class
neighborhoods in (E46.3). The algebraic
isomorphism (E46.21) is consequently a
homeomorphism. This argument permits the
noncompact Kummer completions of number fields.

\emph{The exact total-norm condition.}
Augmentation modulo three gives the exact sequence
of continuous modules
\[
        0\longrightarrow J_A\longrightarrow\Lambda_A
          \xrightarrow{\operatorname{aug}\bmod3}
                     \mathbb F_3(1)\longrightarrow0.
 \tag{E46.22}
\]
The last action is the cyclotomic action; the
translation in (E46.13) has augmentation one.
Since $\mu_3(\mathbb Q)=\{1\}$,
$H^0(G_{\mathbb Q},\mathbb F_3(1))=0$.
Thus $H^1(J_A)$ injects into $H^1(\Lambda_A)$,
with image the kernel of
\[
 H^1(\Lambda_A)\longrightarrow H^1(\mathbb F_3(1))
                     =\mathbb Q^\times/(\mathbb Q^\times)^3.
 \tag{E46.23}
\]
Under (E46.21), this is exactly $u_0$ modulo cubes:
augmentation is the positional map to $n=0$,
which is the total norm from every $E_{A,n}$.
The kernel of
$\widehat{\mathbb Q^\times}_3\to
\mathbb Q^\times/(\mathbb Q^\times)^3$
is $(\widehat{\mathbb Q^\times}_3)^3$.
One can check this without an exactness assumption
on completion: the prime valuations identify
$\widehat{\mathbb Q^\times}_3$ with the sequences
$(v_p)_p\in\prod_p\mathbb Z_3$ tending to zero in
the three-adic sense. Reduction modulo three
has finite support, and a sequence reducing to
zero can be divided by three and still tends
to zero. The rational sign disappears because
cubing is an automorphism of $\{\pm1\}$.
This proves (E46.4).

The injection in (E46.23) is also a topological
identification with its kernel. A cocycle whose
reduction is a coboundary has a unique bounding
element in $\mathbb F_3(1)$, since its invariant
group is zero. Choose a fixed lift of each of
the three possible elements. Subtracting its
coboundary continuously turns the cocycle into
a $J_A$ cocycle. This is an explicit continuous
inverse on the kernel.

\emph{Rational first coordinates and actual ramified sections.}
Choose compatible $3^n$-th roots $a_n$ of $a$ and
compatible $\zeta_{3^n}$. The roots
$a_n^3\zeta_{3^n}^i$ of $X^{3^n}-a^3$ transform by
\[
 g(a_n^3\zeta_{3^n}^i)
     =a_n^3\zeta_{3^n}^{3\kappa_3(a)(g)+\chi_g i}.
\]
Their cube map is the reduction in (E46.2).
This proves precisely the finite-etale algebra
and its transition in (E46.5).

The element $1-X_n$ is invertible in that algebra
because its total norm is $1-a^3\ne0$.
In the rank-three extension given by $Y^3=X_n$,
the determinant of multiplication by $Y$ is $X_n$,
and that of multiplication by $1-Y$,
in the basis $1,Y,Y^2$, is
\[
 \det\begin{pmatrix}
       1&0&-X_n\\-1&1&0\\0&-1&1
      \end{pmatrix}=1-X_n.
 \tag{E46.24}
\]
This determinant proves the norm identity even
when the algebras have several factors. Taking
continuous $\mathbb Z_3$ powers and products proves
compatibility of (E46.6), with
$u_{0,r,z}=a^{3r}(1-a^3)^{3z}$.
The construction (E46.21)--(E46.23) therefore
gives an actual continuous $H$ section class.
Augmentation and (E46.7) give (E46.8).
For rational $a\ne0$ with $a^3\ne1$, the number
$1-a^3$ is neither $1$ nor $-1$: these alternatives
would give $a=0$ or $a^3=2$, respectively.
Its nonzero rational prime-valuation vector
therefore makes $z\mapsto z\kappa_3(1-a^3)$
injective. For $a=2$, the valuations at two and
seven distinguish $r\kappa_3(2)+z\kappa_3(-7)$
for all pairs $(r,z)$. This distinguishes the
stated families through an actual quotient homomorphism.
No identification with an unspecified
rational-point path cocycle is used.

For $a=2$, the factor $X^{3^{n-1}}-2$ is
Eisenstein at two. The complementary roots
are the conjugates of $\beta_n$ in (E46.9).
This element satisfies
$\beta_n^{3^{n-1}}=2\zeta_3$ and contains
$\mathbb Q(\zeta_3)$, since
$\beta_n^{3^{n-1}}/2=\zeta_3$.
The prime two is unramified in
$\mathbb Q(\zeta_3)$; in fact $X^2+X+1$ is
irreducible modulo two. Over that field
$X^{3^{n-1}}-2\zeta_3$ is Eisenstein at its
prime above two. It gives exactly the second
degree in (E46.10). The sum of the two
degrees is $3^n$, proving the full factorization.
Cubing respects both factors and gives the
displayed tower maps. At $n=1$, the total norm
is $v_1N_{L_1/\mathbb Q}w_1$, proving
(E46.11). In particular
$N_{\mathbb Q(\zeta_3)/\mathbb Q}(1-2\zeta_3)=7$
and $1-2=-1$, so the explicit family has
total norm $2^{3r}(-7)^{3z}$, as asserted.

\emph{Finite verification.}
The program \texttt{e\_norm\_sections.py} verifies
the actual affine positional actions, their two
orbits for $a=2$, and the intertwining of all
tested coefficient pushforwards. It implements
(E46.18) in the explicit finite splitting group
of the ninth-root and twenty-seventh-root-of-unity
tower, checks its cocycle identities, and checks
the corestriction-to-norm formula on the nontrivial
Kummer class of $2\zeta_3$ and on the rational
class of $2$. Its $78732$ cocycle identities and
$13122$ inverse-deck comparisons all pass. It also
verifies (E46.24) and the
total-norm and factorization identities at the
listed finite levels. These computations retain
the scalar cyclotomic factor and the direction
of the root-cubing maps. The proofs above give
the asserted all-level comparison.
\end{proof}
